\documentclass{article}
\usepackage{amsmath, amssymb}
\usepackage{cancel}
\usepackage{tikz}
\usepackage{mdframed}
\usepackage{hyperref}
\usetikzlibrary{arrows.meta, decorations.markings}

\begin{document}

\section{Lecture 12: Plan for Class}
Last time:
\begin{enumerate}
    \item Hopf bifurcations
\end{enumerate}

Today:
\begin{enumerate}
    \item Bifurcations of fixed points of maps
\end{enumerate}

\section{Hopf bifurcations}

Here is our example of a system that exhibits a Hopf bifurcation
\begin{align*}
    \dot{x}&= \mu x -\omega y - x \left(x^2 + y^2\right)\\
    \dot{y} &= \omega x + \mu y - y\left(x^2 + y^2\right)
\end{align*}
where the eigenvalues are $\mu\pm i\omega$. Notice that these eigenvalues will cross the imaginary axis as $\mu$ increases. This is the essence of a Hopf bifurcation! Using a change of coordinates to polar coordinates via 
\begin{align*}
    x &= r\cos\theta\\
    y & = r\sin\theta
\end{align*}
to get 
\begin{align*}
    \dot{r}&= r \left(\mu - r^2\right)\\
    \dot{\theta} &= \omega 
\end{align*}
Now, it should be clear that for $\mu\leq 0 $ we get a stable spiral at the origin, and for $\mu > 0$ we get an unstable spiral at the origin and a stable periodic orbit at $r=\sqrt{\mu}$. This is a limit cycle. The is termed a supercritical Hopf bifurcation.

To get a subcritical Hopf bifurcation, we can look at the system
\begin{align*}
    \dot{x}&= \mu x -\omega y + x \left(x^2 + y^2\right)\\
    \dot{y} &= \omega x + \mu y + y\left(x^2 + y^2\right)
\end{align*}
then the change of polar coordinates gets us 
\begin{align*}
    \dot{r}&= r \left(\mu + r^2\right)\\
    \dot{\theta} &= \omega 
\end{align*}
where for $\mu<0$ we have an unstable limit cycle, and for $\mu\geq 0$ we have an unstable spiral at the origin.

\subsection{More generally}

More generally, eigenvalues on the 2nd order center manifold when $\mu=0$ will look like 
\begin{align*}
    \dot{r}&= \left(d\mu + a r^2\right)r \quad = f(r)\\
    \dot{\theta} &= \left(\omega + c \mu + br^2\right)
\end{align*}
where for 
\begin{itemize}
    \item $a<0$ we have a stable limit cycle (supercritical)
    \item $a>0$ we have an unstable periodic orbit (subcritical)
\end{itemize}
To calculate why the value of $a$ determines the stability, we can notice that the equilibrium occurs at $d\mu = -ar_0^2$ and write 
\begin{align*}
    f'(r_0)&= d\mu + 3ar_0^2\\
    &= -ar_0^2 + 3ar_0^2\\
    &= 2ar_0^2
\end{align*}
where it is clear that the sign of $a$ determines the stability of this equilibrium point.

\subsection{Example from last time}

Let's consider the system 
\begin{align*}
    \dot{x}&= -\omega y + \alpha y^2 + \beta x^2 y \\
    \dot{y}&= \omega x - \gamma y^2 + \delta xy - y^3
\end{align*}
where our normal form formula for the value of $a$ gives us 
\[
a=\frac{1}{8}\left(\frac{\gamma}{\omega}\left(\delta - 2\alpha\right)-3\right)
\]
From this expression we can make the following observations
\begin{itemize}
    \item If $\alpha=\gamma=\beta=0$, then there are no quadratic terms in the system we have a stable equilibrium because $a=-\frac{3}{8}$
    \item For large frequency $\omega$ and some quadratic terms, it follows that $a\to -\frac{3}{8}$ and $\omega\to \infty$
    \item For frequency, the system can destabilize. This is apparent by writing $\gamma\left(\delta - 2\alpha\right)>3\omega$ which is an unstable equilibrium
\end{itemize}
This approach is quite useful because linearization tells us nothing and a Lyapunov function would be too complicated to make sense of. Rowley did an example of the Lyapunov function with 
\[
V(x,y) = \frac{1}{2}\left(x^2 + y^2\right)
\]
I leave this as an exercise to the reader.

\section{Bifurcations of fixed points of maps \S 3.5}

For flows $\dot{x}=f(x)$, we got bifurcations when our equilibrium points crossed the imaginary axis. For maps $x\mapsto f(x)$, we will get bifurcations when the \underline{magnitude} of the equilibrium points cross the unit circle.

Much of the same phenomena will transfer over to maps, but we will get a new phenomenon that occurs when the magnitude of an equilibrium point cross $-1$. This is called a flip bifurcation or a period doubling bifurcation.

\subsection{Bifurcation theorems}

Consider the system 
\[
x_{n+1} = f\left(x_n, \mu\right)
\]
where a branch of fixed points $\bar{x}(\mu)$ occurs when 
\[
\bar{x}(\mu) = f\left(\bar{x}(\mu), \mu\right)
\]
Now let's use a Taylor series about the fixed point $x_0=\bar{x}(\mu_0)$ by writing 
\begin{align*}
    f(x_0 + x, \mu_0 + \mu) = f(x_0,\mu_0) &+ \frac{\partial f}{\partial x}(x_0,\mu_0)x + \frac{\partial f}{\partial \mu}(x_0,\mu_0)\mu\\
    &+ \frac{\partial^2 f}{\partial x^2}(x_0,\mu_0)\frac{x^2}{2} + \frac{\partial^2 f}{\partial x \partial \mu}(x_0,\mu_0)x\mu \\
    &+ \frac{\partial^3 f}{\partial x^3}(x_0,\mu_0)\frac{x^3}{3!} + \cdots
\end{align*}
We we have some conditions that we can play around with to get various bifurcations
\begin{itemize}
    \item For all bifurcations: $\frac{\partial f}{\partial x}(x_0,\mu_0)= +1$ (or $-1$, or complex $|\cdot|=1$)
    \item Non-degeneracy: $\frac{\partial^2 f}{\partial x^2}(x_0,\mu_0)\neq 0$
    \item Eigenvalues cross with nonzero speed: $\frac{\partial f}{\partial \mu}\neq0$ but can be complex 
\end{itemize}
Now we can classify the following bifurcations 
\begin{itemize}
    \item Saddle node: $\frac{\partial^2 f}{\partial x^2},\frac{\partial f}{\partial \mu }\neq 0$
    \item Transcritical: $\frac{\partial^2 f}{\partial x^2}\neq 0$ and $\frac{\partial^2 f}{\partial x \partial \mu}\neq 0$, but we could have $\frac{\partial f}{\partial \mu }=0$
    \item Pitchfork: $\frac{\partial^3 f}{\partial x^3}\neq 0$ and $\frac{\partial^2 f}{\partial x \partial \mu}\neq 0$, but we could have $\frac{\partial f}{\partial \mu }=0$ and $\frac{\partial^2 f}{\partial x^2}= 0$
\end{itemize}

\subsubsection{Saddle node example}

We can get a saddle node bifurcation for the following system 
\[
x\mapsto \mu + x \pm x^2
\]
where the equilibrium points occur at
\begin{align*}
    x&= \mu + x \pm x^2\\
    \mu &= \mp x^2
\end{align*}
and the bifurcation occurs when $\mu=0$. 

\subsubsection{Transcritical example}

We can get a transcritical bifurcation in a map at 
\[
x\mapsto \left(1+\mu\right)x\pm x^2
\]
where we see that we have the following conditions for the equilibrium point $(x_0,\mu_0)=(0,0)$
\begin{itemize}
    \item $f_x=1$
    \item $f_{xx}\neq 0$
    \item $f_\mu=0$
    \item $f_{x\mu}\neq 0$
\end{itemize}
The plots are quite informative, but I won't include them.

\subsubsection{Pitchfork example}

We can get a pitchfork bifurcation in a map at 
\[
x\mapsto \left(1+\mu\right)x\pm x^3
\]
where we see that we have the following conditions for the equilibrium point $(x_0,\mu_0)=(0,0)$
\begin{itemize}
    \item $f_x=1$
    \item $f_{xx}=0$
    \item $f_{xxx}\neq 0$
    \item $f_\mu=0$
    \item $f_{x\mu}\neq 0$
\end{itemize}
The plots are quite informative, but I won't include them.

Note that for the positive cubic term $+x^3$, we can compute the equilibrium points as 
\begin{gather*}
    x=\left(1+\mu\right)x + x^3\\
    x= x + \mu x + x^3\\
    0=\mu x + x^3\\
    0=x\left(\mu+x^2\right)
\end{gather*}
where our equilibrium points are $x=0$ and $\mu = - x^2$, where we have a subcritical pitchfork bifurcation.

\section{Flip point bifurcation}

\subsection{Branch of Equilibria}

Consider a branch of equilibria 
\[
\bar{x}(\mu) = f\left( \bar{x}(\mu), \mu \right)
\] 
By differentiating with respect to $\mu$, we have
\begin{gather*}
    \frac{\partial}{\partial \mu}\left[\bar{x}(\mu)\right] = \frac{\partial}{\partial \mu}\left[f\left( \bar{x}(\mu), \mu \right)\right]\\
    \frac{\partial \bar{x}}{\partial \mu} = \frac{\partial f}{\partial x}(\bar{x}(\mu), \mu)\frac{\partial \bar{x}}{\partial \mu} + \frac{\partial f}{\partial \mu}(\bar{x}, \mu)
\end{gather*}
via the chain rule and the product rule. Rearranging some terms, we not have 
\begin{gather*}
    \frac{\partial \bar{x}}{\partial \mu} - \frac{\partial f}{\partial x}(\bar{x}(\mu), \mu)\frac{\partial \bar{x}}{\partial \mu} = \frac{\partial f}{\partial \mu}(\bar{x}, \mu)\\
    \left(1 - \frac{\partial f}{\partial x}(\cdot)\right)\frac{\partial \bar{x}}{\partial \mu} = \frac{\partial f}{\partial \mu}(\bar{x}, \mu)
\end{gather*}
Suppose we want to solve for $\frac{\partial \bar{x}}{\partial \mu}$. We can write 
\[
\frac{\partial \bar{x}}{\partial \mu} = \frac{\frac{\partial f}{\partial \mu}(\bar{x}(\mu), \mu)}{1 - \frac{\partial f}{\partial x}(\bar{x}(\mu), \mu)}
\]
where it follows that we can solve for $\frac{\partial \bar{x}}{\partial \mu}$ as long as 
\[
\frac{\partial f}{\partial x}(x_0, \mu_0) \neq 1
\]
Luckily, at a flip bifurcation, we have 
\[
\frac{\partial f}{\partial x} = -1 \quad \longrightarrow \quad \text{no problem!}
\]
Hence, by the implicit function theorem, it follows that there exists a smooth curve of equilibrium points $\bar{x}(\mu)$.
\begin{mdframed}
    \textbf{Addendum}: The flip bifurcation is "no problem" for the continuation of fixed points — the branch of fixed points doesn't break. What does happen is that the fixed point loses stability and a period-2 orbit is born, but that's a separate calculation.
\end{mdframed}

\subsection{Flip / period doubling bifurcation theorem 3.5.1}

For this theorem, we have three conditions 
\begin{enumerate}
    \item[F0)] $\frac{\partial f}{\partial x}(x_0, \mu_0)=-1$ 
    \item[F1)] Eigenvalue crossing as $\mu$ varies \[ \frac{\partial f}{\partial \mu}\frac{\partial^2 f}{\partial x^2} + 2 \frac{\partial^2 f}{\partial x \partial \mu}=\frac{\partial f}{\partial \mu}\frac{\partial^2 f}{\partial x^2} + \left(1- \frac{\partial f}{\partial x}\right)\frac{\partial^2 f}{\partial x \partial \mu}\neq 0 \]
    \item[F2)] Non-degenerate higher order terms \[ a=\frac{1}{2}\left(\frac{\partial^2 f}{\partial x^2}\right)^2 + \frac{1}{3}\left(\frac{\partial^3 f}{\partial x^3}\right)\neq 0 \]
\end{enumerate}
Why do we need (F1)? Suppose we have a fixed point
\[
\bar{x}(\mu) = f(\bar{x}(\mu),\mu)
\]
and an eigenvalue 
\[
\lambda = \frac{\partial f}{\partial x}(\bar{x}(\mu),\mu)
\]
We want to have $\frac{\partial \lambda}{\partial \mu}\neq 0$, meaning that we cross with non-zero speed. Hence we can write 
\begin{gather*}
    \frac{\partial \lambda}{\partial \mu} = \frac{\partial^2 f}{\partial x^2}\left(\bar{x}(\mu),\mu\right)\frac{\partial \bar{x}}{\mu} + \frac{\partial^2 f}{\partial x \partial \mu}\left(\bar{x}(\mu),\mu\right)\neq 0
\end{gather*}
where we previously derived $\frac{\partial \bar{x}}{\mu}$ as 
\[
\frac{\partial \bar{x}}{\partial \mu} = \frac{\frac{\partial f}{\partial \mu}(\bar{x}(\mu), \mu)}{1 - \frac{\partial f}{\partial x}(\bar{x}(\mu), \mu)}
\]
Plugging in this expression, we have 
\[
\frac{\partial \lambda}{\partial \mu} = \frac{\partial^2 f}{\partial x^2}\left(\bar{x}(\mu),\mu\right)\frac{\frac{\partial f}{\partial \mu}(\bar{x}(\mu), \mu)}{1 - \frac{\partial f}{\partial x}(\bar{x}(\mu), \mu)}+ \frac{\partial^2 f}{\partial x \partial \mu}\left(\bar{x}(\mu),\mu\right)\neq 0
\]
which, by multiplying through by $\left(1-\frac{\partial f}{\partial x}\right)$, is functionally equivalent to 
\[
\frac{\partial f}{\partial \mu}\frac{\partial^2 f}{\partial x^2} + \left(1- \frac{\partial f}{\partial x}\right)\frac{\partial^2 f}{\partial x \partial \mu}=\frac{\partial f}{\partial \mu}\frac{\partial^2 f}{\partial x^2} + 2 \frac{\partial^2 f}{\partial x \partial \mu}\neq 0 
\]
where $\frac{\partial f}{\partial x}=-1\implies \left(1-\frac{\partial f}{\partial x}\right)=2$ as previously noted.

\subsection{Visualizations of period doubling}

For this section, it is necessary to include Rowley's plots.

One things worth mentioning is that period doubling bifurcations can't happen in 2D flows because a trajectory inside a limit cycle would have to intersect the limit cycle to leave. This is topologically not possible.

\subsection{Logistic map}

Let's consider the logistic map 
\[
x\mapsto \mu x \left(1 - x\right)
\]
for values of $\mu>0$. For the rest of the section, Rowley broke out the Jupyter Notebook to numerically look at this system. Let's try to include those plots.

\href{https://doi.org/10.1080/00029890.1975.11994008}{Period three implies chaos}


\end{document}